\section{共享年表之一：东迁、齐霸与晋楚相遇}

\begin{event}{公元前771年}{镐京失守与幽王之死}{通行系年，具体战事叙述有争议}{镐京、骊山一带}\eventanchor{SA-0771-HAO-FALL}{春秋·镐京失守}
西周的终局先表现为王畿失去守备。传统年表把幽王之死、镐京失守安放在这一年；继承冲突、西申关系破裂和西部边防危机相互叠加，使王室不能再把关中旧都作为政治资源的中心。《史记·周本纪》与《竹书纪年》共同保存了这场断裂的年代骨架，《诗经·小雅》所见的危机感则只能提示晚西周的政治环境，不能替代陷落日的军报。

后来的叙事常把申侯、犬戎与褒姒串成一条报应链。它保留了王室联盟破裂的记忆，却不足以让我们复原攻城者的全部身份、路线或宫廷中的每一次选择。能够确定的后果是：西周王畿秩序崩解，王位、宗庙与可动用的人力必须另找落脚处；这正是次年东迁的直接前史。
\end{event}

\begin{sourcenote}[从镐京陷落到东迁：骨架与故事]
前771年幽王死、前770年平王东迁的前后关系在传统年表中稳定；精确月份、继位过程以及烽火、褒姒等完整宫廷故事，则主要来自后出的叙事组织。把这些层次分开，才能既看见关中权力网络的崩裂，也不把兴亡寓言误当作现场记录。
\end{sourcenote}

\begin{event}{公元前770年}{平王东迁洛邑}{可靠纪年}{洛邑}\eventanchor{SA-0770-EAST}{春秋·东迁}
平王迁都洛邑，东周开始。东迁不是把一座旧都搬到东方：宗庙、朝廷和通向诸侯的道路仍可在伊、洛之间重建，但关中原有的仓储、守备与封建网络已不在王室手中。《史记·周本纪》和《竹书纪年》提供迁徙的传统年代线，东周初年的合法性、王畿资源与朝会关系则需要结合《诗经》及历史地理研究谨慎讨论。

王室从关中资源中心退为中原的礼仪与名分中心。此后诸侯仍会借“尊王”来安排会盟、出兵和继承，却越来越少把王命本身当作足以调度军粮和军队的命令。东迁保存了一个可被争用的中心，也把维持这个中心的成本转给了诸侯间不断重组的关系。
\end{event}

东迁后的几十年，传世材料没有给出一条与《春秋》同样连续的列国年表。到前722年鲁国经年开始以后，才较能看见周、郑、宋、卫等国如何把王室职位、边境冲突和会盟互相牵动。前720年，\eventanchor{SA-0720-ZHENG-HOSTAGES}{春秋·周郑交换质子}的传世叙事称周与郑以人质处置互信危机；同年郑又借王师、虢师伐宋，\eventanchor{SA-0720-ZHENG-SONG-ROYAL}{春秋·郑借王师伐宋}。年次和行动可由《春秋·隐公三年》及相关传承安放，质子的姓名、往还言辞与限制效力则主要出自《左传》叙事，不能当作留存的外交文书。

这两件事把东周政治的矛盾放在一处可见的道路上：郑庄公兼有近畿力量与王室卿士位置，王室的军力因而可以进入郑宋边境；可是人质也没有消除双方的猜忌。前719年，宋、陈、蔡、卫联军伐郑，鲁翚会师，\eventanchor{SA-0719-ANTI-ZHENG}{春秋·宋陈蔡卫伐郑}。《春秋·隐公四年》可确认参战者，却不能告诉我们各军究竟出动多少兵车，或谁承担了供给。卫州吁以外战谋求国内承认、宋郑旧怨未解，是解释这次联军的条件，不是由相邻年份自动推出的单一原因。

前715年，宋、齐、卫盟于瓦屋，\eventanchor{SA-0715-SONG-WEI-ALLIANCE}{春秋·宋齐卫瓦屋之盟}；前710年，诸侯会于稷处理宋国之乱，\eventanchor{SA-0710-JI-SETTLE-SONG}{春秋·稷会调停宋乱}。经文留下了会盟与会见的年份，后出的《左传》才展开宋督弑君、继承失序及调停的叙事。会盟因此不是一张已经保存下来的国际条约：它让各国暂时把宋的内乱纳入共同处置，却没有消除宋国大夫政治，也不能证明参会者承担同额兵役、贡赋或长期从属。

\begin{event}{公元前707年}{繻葛之战}{可靠纪年}{繻葛}\eventanchor{SA-0707-XUGE}{春秋·繻葛之战}
周郑之间原有的猜忌终于转为公开军事冲突。此前蔡、卫、陈随王师伐郑，\eventanchor{SA-0707-WANG-ALLIES-ATTACK-ZHENG}{春秋·王师诸侯伐郑}，使周王室与郑的争田、卿士权力和近畿资源之争取得了联盟背景。《春秋·桓公五年》可以确认王师与诸侯伐郑，却不能据此断言周王亲自统率所有军队，亦无从核定兵额、行军次序和各国的实际投入。

同年郑庄公与周桓王战于繻葛，王师败绩。经文的“王师败绩于繻葛”足以固定结果；《左传》再加入桓王中箭、郑庄公不追击以及礼仪往还等战术和人物细节，《史记·郑世家》又从汉代把它纳入郑国兴衰线。战场地望与周郑的质子、麦禾争端仍是历史地理和文本层次研究的问题，不能把传世故事拼成逐时战报。

这一仗没有废除周天子的名义，却让诸侯看见名义不能自动兑换为军事实力。王室仍能召集盟友，郑也仍须在近畿处理王命；但掌握地方军政资源的诸侯和卿士已经能够拒绝王师，东迁后王室依赖诸侯的代价由此更清楚地显现。
\end{event}

\begin{sourcenote}[繻葛：经文与叙事的距离]
经文足以确认王师失败，并保存战争的鲁历年位置；《左传》提供箭伤、礼仪和人物态度，让我们看到后世如何解释周郑关系，却带有鲜明的政治伦理章法。可以用它追问冲突如何被理解，不把其中每句对白都当成战场录音。
\end{sourcenote}

\begin{event}{公元前685年}{齐桓公即位，管仲入齐}{可靠纪年}{齐}\eventanchor{SA-0685-QIHUAN}{春秋·齐桓公即位}
齐襄公被杀、公孙无知短暂据位后，公子小白先于公子纠进入齐国即位，是为齐桓公。相关《春秋》经年及《左传·庄公九年》给出政变与继承的年代骨架；《左传》把小白、纠与管仲的竞争组织成连续故事，《国语·齐语》保存“修政以服诸侯”的齐国政治记忆，《史记·齐太公世家》则从汉代回望，把它收束为霸业的开端。它们相互照亮，却都不是宫门内逐时留下的记录。

齐能在内乱后重建中心，所凭借的不只是国君的个人声望。滨海资源、内政整合和能够召集邻国的道路，才让“尊王攘夷”的语言有了可以兑现的组织条件；管仲在这一政治转折中的位置可靠，具体政策却须分开衡量。
\end{event}

\begin{debate}[管仲改革能不能全信《管子》]
《国语·齐语》和《左传》可帮助理解齐的组织能力；《管子》则是战国至汉代逐层形成的复合文本，不能把其中精密的盐铁、货币和经济理论直接倒灌成前685年管仲的原始政策。管仲确实是齐霸的重要人物，但“具体做了哪一条现代意义上的经济改革”必须保持节制。
\end{debate}

齐的优势并非只由大国彼此认可。前678年齐会诸侯于幽，许、滑、滕等小国与宋、陈、卫、郑一起进入更宽的关系网，\eventanchor{SA-0678-YOU-ALLIANCE}{春秋·幽会诸侯}。《春秋·庄公十六年》可确认参会国别；《左传》的铺陈以及会盟地望研究，才帮助我们追问它怎样把齐对宋的短期武力转为多层联盟。小国加入会盟可能获得保护，也可能承担道路、粮秣和外交上的压力；仅凭参会名单，不能把它们写成齐的固定属国，或推定同额的兵役与贡赋。

\begin{event}{公元前632年}{城濮之战}{可靠纪年}{城濮}\eventanchor{SA-0632-CHENGPU}{春秋·城濮之战}
晋、楚在城濮交战，晋文公一方获胜。《春秋·僖公二十八年》确认战事、参战者与楚师败绩，构成最小的经年骨架；晋能把流亡时期结成的关系转为联盟，又须穿过曹、卫等地取得补给和通道，楚方则面对北上关系能否维持的政治条件。这些都是理解胜负的结构性因素，不能缩减为一位君主的临阵谋略。

《左传》写出的“退三舍”，把晋文公流亡时对楚的承诺变成战前退让的道德高潮；《国语·晋语》以国别说辞保留晋的政治记忆，《史记·晋世家》《楚世家》再把战争置入汉代的列国兴亡线。它们说明后人如何赋予这场战役意义，却不能替代战场日志。城濮之后，晋成为中原联盟的关键组织者；楚则显示南方大国已经无法被排除在“天下”之外。
\end{event}

胜利随后须被组织成诸侯能够承认的次序。城濮之后，晋文公会诸侯于践土并受周王策命，\eventanchor{SA-0632-JIANTU}{春秋·践土之会}。《春秋·僖公二十八年》可确认会盟与其年次，《左传》《国语》才给出王命措辞和礼仪的完整面貌，《史记·晋世家》则把它纳入晋文公的霸业。周襄王需要强藩维持王室名分，晋也需要把军事胜利转换为可协调诸侯的语言；两者的需要在践土相接，却不等于周王把一个现代法定职位“册封”给晋。

践土的效力依然取决于盟国履约、军队供给与下一次危机中的选择。它没有让晋直接征税、任命或裁决所有诸侯，也没有使楚的实力消失；它表明的是军事资源、会盟程序和王室名分在某一时刻可以被合拢。此后的竞争仍要在道路、粮秣、国内卿族和不断变动的盟友之间重新支付成本。

\begin{debate}[“退三舍”与践土的名分]
退三舍只见于《左传》的完整叙事，不在《春秋》极简经文中。它可能保留了某种外交或撤退行为的记忆，但“守信而获胜”的结构高度符合《左传》的伦理叙事。相应地，践土会盟和周王名分语言能够说明晋的协调资格上升，却不能直接证明诸侯从此长期服从，或存在一项稳定、制度化的“霸主”职位。
\end{debate}
